% !TeX root = tcolorbox.tex
% include file of tcolorbox.tex (manual of the LaTeX package tcolorbox)
\clearpage
\section{程序库 \mylib{skins}——皮肤目录}\label{sec:skincatalog}%
\tcbset{external/prefix=external/skincatalog_}%
程序库 \mylib{skins} 提供了一组可供选择的皮肤目录，说明如下。
该程序库需通过宏包选项或在导言区中加载：
\begin{dispListing}
\tcbuselibrary{skins}
\end{dispListing}

有关 \mylib{skins} 程序库其他选项的说明，见 \zvref{sec:skins}。

\begin{itemize}
\item 原则上，皮肤通过为 \refKey{/tcb/skin} 选择一个值来应用，例如 \docValue*{enhanced}。
  由于可断开文本框的各个片段应具有不同外观，
  因此为可断开文本框分别设置了独立的皮肤，另见 \zvref{subsec:breaksequence}。
  由同一基础皮肤派生出来的、用于可断开文本框的皮肤，下文称为一个皮肤家族。
\item 除了为 \refKey{/tcb/skin} 赋值之外，还可使用同名的选项，它们是快捷方式，
  有时还会一并改变某些几何或样式设置。这些才是供普通用户使用的选项。
  通常，选用下列选项之一即可选定一种皮肤：
  \begin{itemize}
  \item \refKey{/tcb/standard}
  \item \refKey{/tcb/standard jigsaw}
  \item \refKey{/tcb/enhanced}
  \item \refKey{/tcb/enhanced jigsaw}
  \item \refKey{/tcb/enhanced standard}
  \item \refKey{/tcb/enhanced standard jigsaw}
  \item \refKey{/tcb/bicolor}
  \item \refKey{/tcb/tile}
  \item \refKey{/tcb/beamer}
  \item \refKey{/tcb/widget}
  \item \refKey{/tcb/empty}
  \item \refKey{/tcb/spartan}
  \item \refKey{/tcb/draft}
  \end{itemize}
  此外，还有一些特殊用途：
  \begin{itemize}
  \item \refKey{/tcb/marker}
  \item \refKey{/tcb/blank}
  \item \refKey{/tcb/blanker}
  \item \refKey{/tcb/blankest}
  \end{itemize}
\end{itemize}



\clearpage

下面的示例使用辅助宏 \docAuxCommand{skinExampleSet} 来展示皮肤效果。
请注意，\docAuxCommand{skinExampleSet} 并非宏包的一部分，
而只是为本文档定义的。

\begin{dispListing}
\NewDocumentCommand{\skinExampleSet}{m}{%
  \begin{tcbraster}[raster equal height,raster columns=3,
      colback=LightGreen,colframe=DarkGreen,colbacktitle=LimeGreen!75!DarkGreen,
      #1,
      left=1mm,right=1mm,top=1mm,bottom=1mm,middle=1mm,
      sidebyside gap=4mm]
    \begin{tcolorbox}
      This is my content.
    \end{tcolorbox}
    \begin{tcolorbox}
      This is my content.
      \tcblower
      More content.
    \end{tcolorbox}
    \begin{tcolorbox}[sidebyside]
      My content.
      \tcblower
      More content.
    \end{tcolorbox}
    \begin{tcolorbox}[adjusted title=My title]
      This is my content.
    \end{tcolorbox}
    \begin{tcolorbox}[adjusted title=My title]
      This is my content.
      \tcblower
      More content.
    \end{tcolorbox}
    \begin{tcolorbox}[adjusted title=My title,sidebyside]
      My content.
      \tcblower
      More content.
    \end{tcolorbox}
  \end{tcbraster}
}
\end{dispListing}
\tcbusetemp


\clearpage
\tcbset{skintable/.style={colframe=red!50!yellow!50!black,
  colback=red!50!yellow!5!white,coltitle=red!50!yellow!3!white,
  fonttitle=\bfseries,before=\par\smallskip,
  title=皮肤 \enquote{\texttt{#1}} 的环境与引擎}}

\subsection{皮肤家族 \enquote{standard}}\label{subsec:skinstandard}
\begin{marker}请注意，选项键 \refKey{/tcb/frame style}、
  \refKey{/tcb/interior style}、
  \refKey{/tcb/segmentation style} 和
  \refKey{/tcb/title style} 不适用于 standard 皮肤。
  此外，水印（见 \zvref{subsec:watermarks}）
  也不能用于 standard 皮肤。
\end{marker}

\begin{docSkin}{standard}
  这是核心宏包的标准皮肤。所有绘制引擎均设为 |standard| 类型。
  绘制基于 \refPkg{pgf} 命令，不需要 \refPkg{tikz} 宏包。
\begin{tcolorbox}[skintable=standard]
  \begin{tabbing}
    \refKey{/tcb/interior titled engine}: \=\kill
    \refKey{/tcb/frame engine}:           \> |standard|\\
    \refKey{/tcb/interior titled engine}: \> |standard|\\ 
    \refKey{/tcb/interior engine}:        \> |standard|\\
    \refKey{/tcb/segmentation engine}:    \> |standard|\\
    \refKey{/tcb/title engine}:           \> |standard|
  \end{tabbing}
\end{tcolorbox}
\end{docSkin}

\begin{docTcbKey}{standard}{}{样式，无取值}
  这是设置 |skin=standard| 的简写。
\end{docTcbKey}

\begin{dispExample}
\skinExampleSet{standard}
\end{dispExample}

\clearpage

\begin{docSkin}{standard jigsaw}
  这是核心宏包的标准拼图皮肤。它与 \refSkin{standard} 皮肤的区别在于其边框引擎，
  见 \zvref{subsec:skinjigsaw}。
\begin{tcolorbox}[skintable=standard jigsaw]
  \begin{tabbing}
    \refKey{/tcb/interior titled engine}: \=\kill
    \refKey{/tcb/frame engine}:           \> |standardjigsaw|\\
    \refKey{/tcb/interior titled engine}: \> |standard|\\ 
    \refKey{/tcb/interior engine}:        \> |standard|\\
    \refKey{/tcb/segmentation engine}:    \> |standard|\\
    \refKey{/tcb/title engine}:           \> |standard|
  \end{tabbing}
\end{tcolorbox}
\end{docSkin}

\begin{docTcbKey}{standard jigsaw}{}{样式，无取值}
  这是设置 |skin=standard jigsaw| 的简写。
\end{docTcbKey}

\begin{dispExample*}{segmentation style={preaction={fill=white},pattern=checkerboard,pattern color=gray!40}}
\skinExampleSet{standard jigsaw,
  opacityframe=0.5,opacityback=0.5,opacitybacktitle=0.5,
}
\end{dispExample*}


\clearpage
\subsection{皮肤家族 \enquote{enhanced}}
\begin{marker}
如果你喜欢 |tcolorbox| 的标准外观，但又希望拥有一些\enquote{增强}特性，
那么 |enhanced| 皮肤正是你想要的。
\end{marker}

\begin{docSkin}{enhanced}
  该皮肤把核心宏包的绘制命令转换为 \tikzname\ 路径命令。因此，它允许对这些路径
  使用所有 \tikzname\ 高级选项，相比 \refSkin{standard} 皮肤更加灵活。
  代价是编译时间会稍长。这些 \tikzname\ 路径选项可通过下列选项键给出：
  \refKey{/tcb/frame style}、
  \refKey{/tcb/interior style}、
  \refKey{/tcb/segmentation style} 和
  \refKey{/tcb/title style}。
\begin{tcolorbox}[skintable=enhanced]
  \begin{tabbing}
    \refKey{/tcb/interior titled engine}: \=\kill
    \refKey{/tcb/frame engine}:           \> |path|\\
    \refKey{/tcb/interior titled engine}: \> |path|\\ 
    \refKey{/tcb/interior engine}:        \> |path|\\
    \refKey{/tcb/segmentation engine}:    \> |path|\\
    \refKey{/tcb/title engine}:           \> |path|
  \end{tabbing}
\end{tcolorbox}
\end{docSkin}


\begin{docTcbKey}{enhanced}{}{样式，无取值}
  这是设置 |skin=enhanced| 的简写。
\end{docTcbKey}

\begin{dispExample}
\skinExampleSet{enhanced}
\end{dispExample}

\clearpage

\begin{dispExample}
% \usetikzlibrary{shadings}         % preamble
\tcbset{skin=enhanced,fonttitle=\bfseries,
  frame style={upper left=blue,upper right=red,lower left=yellow,lower right=green},
  interior style={white,opacity=0.5},
  segmentation style={black,solid,opacity=0.2,line width=1pt}}

\begin{tcolorbox}[title=Nice box in rainbow colors]
  With the \enquote{enhanced} skin, it is quite easy to produce fancy looking effects.
  \tcblower
  Note that this is still a \texttt{tcolorbox}.
\end{tcolorbox}
\end{dispExample}


\begin{dispExample}
% \usetikzlibrary{decorations.pathmorphing} % preamble
\tcbset{skin=enhanced,fonttitle=\bfseries,boxrule=1mm,
  frame style={draw=FireBrick,fill=Salmon},drop fuzzy shadow,
  interior style={draw=FireBrick,top color=Salmon!10,bottom color=Salmon!20},
  segmentation style={draw=FireBrick,solid,decorate,
        decoration={coil,aspect=0,segment length=10.1mm}}}

\begin{tcblisting}{title=A listing box with shadow and some specials}
Of course, skins can be used for listings also.
\begin{equation}
  \int\limits_1^2 \frac{1}{x}~dx = \ln(2).
\end{equation}
\end{tcblisting}
\end{dispExample}


\clearpage


\begin{docTcbKey}{enhanced standard}{}{样式，无取值}
  本样式为所有文本框（包括不可断开的文本框以及\emph{断开序列}中的每一段）
  统一设置未断开框所用的 \refSkin{enhanced} 皮肤。因此，它是
  \refSkin{standard} 皮肤的另一种变体，见 \zvref{subsec:breaksequence}。
\end{docTcbKey}


\begin{docTcbKey}{blank}{}{样式，初始未设置}
  本样式基于 \refSkin{enhanced} 皮肤。所有绘制操作均被隐藏，
  所有边距均设为 |0pt|。要关闭绘制引擎，见 \refKey{/tcb/blanker}。
\begin{dispExample}
\begin{tcolorbox}[blank,watermark text=A blank box]
\lipsum[1]
\end{tcolorbox}
\end{dispExample}
\end{docTcbKey}

\clearpage
\begin{docCommand}{tcbline}{}
  有时候，一条线就是一条线。使用 \refCom{tcblower} 可以把文本框内容分成两个功能单元。
  |\tcbline| 只画一条线，其外观类似于上下部分之间的分隔线。
  此外，|\tcbline| 可以使用不止一次。
  |\tcbline| 始终使用 |path| 绘制引擎，因此可以应用 \refKey{/tcb/segmentation style}。

\begin{dispExample}
\tcbset{enhanced,colframe=blue!50!black,colback=white}

\begin{tcolorbox}[colupper=red!50!black,collower=green!50!black]
\lipsum[1]
\tcbline
\lipsum[2]
\tcblower
\lipsum[3]
\tcbline
\lipsum[4]
\end{tcolorbox}
\end{dispExample}
\end{docCommand}

\begin{docCommand}{tcbline*}{}
  与 \refCom{tcbline} 相同，但在可断开文本框中，若发生页面/框断开，
  \refCom{tcbline*} 会被移除；在文本框末尾也会被移除。
\end{docCommand}

\clearpage
\begin{docSkin}{enhancedfirst}
这是 \refSkin{enhanced} 的一个变体，用作 \refSkin{enhanced} 断开序列的\emph{首}段。
不过，该皮肤也可以独立使用。
\begin{tcolorbox}[skintable=enhancedfirst]
  \begin{tabbing}
    \refKey{/tcb/interior titled engine}: \=\kill
    \refKey{/tcb/frame engine}:           \> |pathfirst|\\
    \refKey{/tcb/interior titled engine}: \> |pathfirst|\\ 
    \refKey{/tcb/interior engine}:        \> |pathfirst|\\
    \refKey{/tcb/segmentation engine}:    \> |path|\\
    \refKey{/tcb/title engine}:           \> |pathfirst|
  \end{tabbing}
\end{tcolorbox}
\end{docSkin}


\begin{dispExample}
\skinExampleSet{skin=enhancedfirst}
\end{dispExample}

\medskip

%\clearpage
\begin{docSkin}{enhancedmiddle}
这是 \refSkin{enhanced} 的一个变体，用作 \refSkin{enhanced} 断开序列的\emph{中间}段。
不过，该皮肤也可以独立使用。
\begin{tcolorbox}[skintable=enhancedmiddle]
  \begin{tabbing}
    \refKey{/tcb/interior titled engine}: \=\kill
    \refKey{/tcb/frame engine}:           \> |pathmiddle|\\
    \refKey{/tcb/interior titled engine}: \> |pathmiddle|\\ 
    \refKey{/tcb/interior engine}:        \> |pathmiddle|\\
    \refKey{/tcb/segmentation engine}:    \> |path|\\
    \refKey{/tcb/title engine}:           \> |pathmiddle|
  \end{tabbing}
\end{tcolorbox}
\end{docSkin}


\begin{dispExample}
\skinExampleSet{skin=enhancedmiddle}
\end{dispExample}


\clearpage
\begin{docSkin}{enhancedlast}
这是 \refSkin{enhanced} 的一个变体，用作 \refSkin{enhanced} 断开序列的\emph{末}段。
不过，该皮肤也可以独立使用。
\begin{tcolorbox}[skintable=enhancedlast]
  \begin{tabbing}
    \refKey{/tcb/interior titled engine}: \=\kill
    \refKey{/tcb/frame engine}:           \> |pathlast|\\
    \refKey{/tcb/interior titled engine}: \> |pathlast|\\ 
    \refKey{/tcb/interior engine}:        \> |pathlast|\\
    \refKey{/tcb/segmentation engine}:    \> |path|\\
    \refKey{/tcb/title engine}:           \> |pathlast|
  \end{tabbing}
\end{tcolorbox}
\end{docSkin}

\begin{dispExample}
\skinExampleSet{skin=enhancedlast}
\end{dispExample}

\clearpage

\begin{docSkin}{enhanced jigsaw}
  这是 \refSkin{enhanced} 皮肤的拼图变体。区别在于其边框引擎，
  见 \zvref{subsec:skinjigsaw}。
\begin{tcolorbox}[skintable=enhanced jigsaw]
  \begin{tabbing}
    \refKey{/tcb/interior titled engine}: \=\kill
    \refKey{/tcb/frame engine}:           \> |pathjigsaw|\\
    \refKey{/tcb/interior titled engine}: \> |path|\\ 
    \refKey{/tcb/interior engine}:        \> |path|\\
    \refKey{/tcb/segmentation engine}:    \> |path|\\
    \refKey{/tcb/title engine}:           \> |path|
  \end{tabbing}
\end{tcolorbox}
\end{docSkin}

\begin{docTcbKey}{enhanced jigsaw}{}{样式，无取值}
  这是设置 |skin=enhanced jigsaw| 的简写。
\end{docTcbKey}


\begin{dispExample*}{segmentation style={preaction={fill=white},pattern=checkerboard,pattern color=gray!40}}
\skinExampleSet{enhanced jigsaw,
  opacityframe=0.5,opacityback=0.5,opacitybacktitle=0.5,
}
\end{dispExample*}


\begin{docTcbKey}[][doc new=2017-07-01]{enhanced standard jigsaw}{}{样式，无取值}
  本样式为所有文本框（包括不可断开的文本框以及\emph{断开序列}中的每一段）
  统一设置未断开框所用的 \refSkin{enhanced jigsaw} 皮肤。因此，它是
  \refSkin{standard} 皮肤的另一种变体，见 \zvref{subsec:breaksequence}。
\end{docTcbKey}


\clearpage
\begin{docSkin}{enhancedfirst jigsaw}
  这是 \refSkin{enhancedfirst} 皮肤的拼图变体。区别在于其边框引擎，
  见 \zvref{subsec:skinjigsaw}。
\begin{tcolorbox}[skintable=enhancedfirst jigsaw]
  \begin{tabbing}
    \refKey{/tcb/interior titled engine}: \=\kill
    \refKey{/tcb/frame engine}:           \> |pathfirstjigsaw|\\
    \refKey{/tcb/interior titled engine}: \> |pathfirst|\\ 
    \refKey{/tcb/interior engine}:        \> |pathfirst|\\
    \refKey{/tcb/segmentation engine}:    \> |path|\\
    \refKey{/tcb/title engine}:           \> |pathfirst|
  \end{tabbing}
\end{tcolorbox}
\end{docSkin}


\begin{dispExample*}{segmentation style={preaction={fill=white},pattern=checkerboard,pattern color=gray!40}}
\skinExampleSet{skin=enhancedfirst jigsaw,
  opacityframe=0.5,opacityback=0.5,opacitybacktitle=0.5,
}
\end{dispExample*}


\clearpage
\begin{docSkin}{enhancedmiddle jigsaw}
  这是 \refSkin{enhancedmiddle} 皮肤的拼图变体。区别在于其边框引擎，
  见 \zvref{subsec:skinjigsaw}。
\begin{tcolorbox}[skintable=enhancedmiddle jigsaw]
  \begin{tabbing}
    \refKey{/tcb/interior titled engine}: \=\kill
    \refKey{/tcb/frame engine}:           \> |pathmiddlejigsaw|\\
    \refKey{/tcb/interior titled engine}: \> |pathmiddle|\\ 
    \refKey{/tcb/interior engine}:        \> |pathmiddle|\\
    \refKey{/tcb/segmentation engine}:    \> |path|\\
    \refKey{/tcb/title engine}:           \> |pathmiddle|
  \end{tabbing}
\end{tcolorbox}
\end{docSkin}


\begin{dispExample*}{segmentation style={preaction={fill=white},pattern=checkerboard,pattern color=gray!40}}
\skinExampleSet{skin=enhancedmiddle jigsaw,
  opacityframe=0.5,opacityback=0.5,opacitybacktitle=0.5,
}
\end{dispExample*}


\begin{docTcbKey}{marker}{}{样式，无取值}
  本样式基于 \refSkin{enhancedmiddle jigsaw} 皮肤。
  它旨在像荧光笔一样用作视觉标记。
\begin{dispExample}
\begin{tcolorbox}[marker]
\lipsum[2]
\end{tcolorbox}
\end{dispExample}
\end{docTcbKey}

\clearpage

\begin{dispListing*}{before upper={本示例展示如何基于 \refSkin{enhancedmiddle} 创建若干
  \emph{文本标记} 环境。\par\medskip}}
\tcbset{textmarker/.style={%
    skin=enhancedmiddle jigsaw,breakable,parbox=false,
    boxrule=0mm,leftrule=5mm,rightrule=5mm,boxsep=0mm,arc=0mm,outer arc=0mm,
    left=3mm,right=3mm,top=1mm,bottom=1mm,toptitle=1mm,bottomtitle=1mm,oversize}}

\newtcolorbox{yellow}{textmarker,colback=yellow!5!white,colframe=yellow}
\newtcolorbox{orange}{textmarker,colback=DarkOrange!5!white,
                        colframe=DarkOrange!75!yellow}
\newtcolorbox{red}{textmarker,colback=red!5!white,colframe=red}
\newtcolorbox{blue}{textmarker,colback=DeepSkyBlue!5!white,colframe=DeepSkyBlue}
\newtcolorbox{green}{textmarker,colback=Chartreuse!5!white,colframe=Chartreuse}
\newtcolorbox{rainbow}{textmarker,interior hidden,
  frame style={top color=blue,bottom color=red,middle color=green}}

\begin{yellow}
  \lipsum[1-3]
\end{yellow}

\begin{orange}
  \lipsum[4]
\end{orange}

\begin{red}
  \lipsum[5]
\end{red}

\begin{green}
  \lipsum[6]
\end{green}

\begin{blue}
  \lipsum[7]
\end{blue}

\begin{rainbow}
  \lipsum[8]
\end{rainbow}
\end{dispListing*}
{\tcbusetemp}


\clearpage
\begin{docSkin}{enhancedlast jigsaw}
  这是 \refSkin{enhancedlast} 皮肤的拼图变体。区别在于其边框引擎，
  见 \zvref{subsec:skinjigsaw}。
\begin{tcolorbox}[skintable=enhancedlast]
  \begin{tabbing}
    \refKey{/tcb/interior titled engine}: \=\kill
    \refKey{/tcb/frame engine}:           \> |pathlastjigsaw|\\
    \refKey{/tcb/interior titled engine}: \> |pathlast|\\ 
    \refKey{/tcb/interior engine}:        \> |pathlast|\\
    \refKey{/tcb/segmentation engine}:    \> |path|\\
    \refKey{/tcb/title engine}:           \> |pathlast|
  \end{tabbing}
\end{tcolorbox}
\end{docSkin}


\begin{dispExample*}{segmentation style={preaction={fill=white},pattern=checkerboard,pattern color=gray!40}}
\skinExampleSet{skin=enhancedlast jigsaw,
  opacityframe=0.5,opacityback=0.5,opacitybacktitle=0.5,
}
\end{dispExample*}



\clearpage
\subsection{皮肤家族 \enquote{bicolor}}
\begin{docSkin}{bicolor}
  该皮肤与 \refSkin{standard} 和 \refSkin{enhanced} 皮肤十分相似。
  但它不用分隔线，而是用不同颜色填充或用不同样式绘制文本框的可选下部。
\begin{tcolorbox}[skintable=bicolor]
  \begin{tabbing}
    \refKey{/tcb/interior titled engine}: \=\kill
    \refKey{/tcb/frame engine}:           \> |path|\\
    \refKey{/tcb/interior titled engine}: \> \emph{special}\\ 
    \refKey{/tcb/interior engine}:        \> \emph{special}\\
    \refKey{/tcb/segmentation engine}:    \> \emph{special}\\
    \refKey{/tcb/title engine}:           \> |path|
  \end{tabbing}
\end{tcolorbox}
  \begin{itemize}
  \item 该皮肤最基本的用法是用 \refKey{/tcb/colbacklower} 设置下部的背景色，
    其余选项与 \refSkin{standard} 皮肤相同。
\begin{dispExample}
\begin{tcolorbox}[skin=bicolor,title=The title,
    colframe=FireBrick!75!black,colback=Salmon!50!white,colbacklower=Salmon]
  The upper part.
  \tcblower
  The lower part.
\end{tcolorbox}
\end{dispExample}
  \item 更高级的用法是像 \refSkin{enhanced} 皮肤那样应用 \refKey{/tcb/frame style}
    和 \refKey{/tcb/interior style}。也可以使用 \refKey{/tcb/segmentation style}，
    但它会作用于整个下部。
\begin{dispExample}
\begin{tcolorbox}[skin=bicolor,title=The title,
    frame style={top color=FireBrick,
                 bottom color=FireBrick!15!white,draw=black},
    interior style={left color=Salmon,right color=Salmon!50!white},
    segmentation style={right color=Salmon,left color=Salmon!50!white}]
  The upper part.
  \tcblower
  The lower part.
\end{tcolorbox}
\end{dispExample}
  \end{itemize}
\end{docSkin}

\clearpage

\begin{docTcbKey}{bicolor}{}{样式，无取值}
  这是设置 |skin=bicolor| 的简写。
\end{docTcbKey}


\begin{dispExample}
\skinExampleSet{bicolor,
  colbacklower=LimeGreen!75!LightGreen,
}
\end{dispExample}

\clearpage


\begin{marker}
  下列选项 \refKey{/tcb/colbacklower} 和 \refKey{/tcb/opacitybacklower}
  在 \refKey{/tcb/segmentation style} 之前执行，也就是说，
  \refKey{/tcb/segmentation style} 会覆盖它们。
\end{marker}

\begin{docTcbKey}{colbacklower}{=\meta{color}}{无默认值，初始 \texttt{black!15!white}}
  设置下部的背景 \meta{颜色}。是否使用该值取决于所用皮肤。
\end{docTcbKey}

\begin{dispExample}
\tcbset{gitexample/.style={listing and comment,comment={#1},
  skin=bicolor,boxrule=1mm,fonttitle=\bfseries,coltitle=black,
  frame style={draw=black,left color=Gold,right color=Goldenrod!50!Gold},
  colback=black,colbacklower=Goldenrod!75!Gold,
  colupper=white,collower=black,
  listing options={language={bash},aboveskip=0pt,belowskip=0pt,nolol,
  basicstyle=\ttfamily\bfseries,extendedchars=true}}}

\begin{tcblisting}{title={Snapshot of the staging area},
  gitexample={The option `-a' automatically stages all tracked and modified
              files before the commit.\par
              This can be combined with the message option `-m'
              as seen in the third line.}}
git commit
git commit -a
git commit -am 'changes to my example'
\end{tcblisting}
\end{dispExample}

\smallskip

\begin{docTcbKey}[][doc new=2021-05-21]{opacitybacklower}{=\meta{fraction}}{无默认值，初始 \texttt{1.0}}
  将下部的背景不透明度设为所给 \meta{比例}。
  是否使用该值取决于所用皮肤。
\end{docTcbKey}

\begin{dispExample}
\begin{tcolorbox}[bicolor,
  frame style={preaction={fill=blue!50!black},
    pattern=checkerboard,pattern color=blue!50!gray},
  fonttitle=\bfseries, overlaplower=0mm,
  colback=blue!10, colbacklower=white, opacitybacklower=0.65,
  title={Example for a semilucent lower part}]
This is the upper part.
\tcblower
And that is the lower part.
\end{tcolorbox}
\end{dispExample}

\clearpage

\begin{docTcbKey}[][doc new=2022-01-11]{overlaplower}{=\meta{length}}{无默认值，初始 \texttt{0.1mm}}
  对于皮肤家族 \enquote{bicolor}、\enquote{tile} 和 \enquote{beamer}，
  下部背景的绘制方式与上部背景不同。若上下部背景之间的距离为 |0pt|，
  某些预览器会在上下部之间显示一条细的边框色线。为避免这一瑕疵，
  下部会以 \meta{长度} 的重叠量绘制在上部之上。\par
  该值可根据特殊需要调整。例如，半透明下部最好使用 |0pt|，
  见 \refKey{/tcb/opacitybacklower}。要使用更大的值，另见 \refCom{tcboverlaplower}。
\end{docTcbKey}


\begin{docCommand}[doc new=2022-01-11]{tcboverlaplower}{}
  该宏保存由 \refKey{/tcb/overlaplower} 设置的长度值。
  可用于上下部分隔处的精细定位，应视为\emph{只读}。

\begin{dispExample}
\begin{tcolorbox}[bicolor, sharp corners,
    colframe=blue!50!black, colback=blue!10, colbacklower=red!10,
    top=5mm, bottom=2mm, middle=3.5mm, overlaplower=1.5mm,
    underlay={
      \node[minimum width=1cm,minimum height=0.5cm,outer sep=auto,
        anchor=north east,fill=white] at (interior.north east)
        {\itshape\small upper};
      \node[minimum width=1cm,minimum height=0.5cm,outer sep=auto,
        anchor=north east,fill=white]
        at ([yshift=\tcboverlaplower]segmentation.east)
        {\itshape\small lower};
    }
  ]
This is the upper part.
\tcblower
And that is the lower part.
\end{tcolorbox}
\end{dispExample}

\end{docCommand}



\clearpage

\begin{docSkin}{bicolorfirst}
这是 \refSkin{bicolor} 的一个变体，用作 \refSkin{bicolor} 断开序列的\emph{首}段。
不过，该皮肤也可以独立使用。
\begin{tcolorbox}[skintable=bicolorfirst]
  \begin{tabbing}
    \refKey{/tcb/interior titled engine}: \=\kill
    \refKey{/tcb/frame engine}:           \> |pathfirst|\\
    \refKey{/tcb/interior titled engine}: \> \emph{special}\\ 
    \refKey{/tcb/interior engine}:        \> \emph{special}\\
    \refKey{/tcb/segmentation engine}:    \> \emph{special}\\
    \refKey{/tcb/title engine}:           \> |pathfirst|
  \end{tabbing}
\end{tcolorbox}
\end{docSkin}

\begin{dispExample}
\skinExampleSet{skin=bicolorfirst,
  colbacklower=LimeGreen!75!LightGreen,
}
\end{dispExample}


\clearpage

\begin{docSkin}{bicolormiddle}
这是 \refSkin{bicolor} 的一个变体，用作 \refSkin{bicolor} 断开序列的\emph{中间}段。
不过，该皮肤也可以独立使用。
\begin{tcolorbox}[skintable=bicolormiddle]
  \begin{tabbing}
    \refKey{/tcb/interior titled engine}: \=\kill
    \refKey{/tcb/frame engine}:           \> |pathmiddle|\\
    \refKey{/tcb/interior titled engine}: \> \emph{special}\\ 
    \refKey{/tcb/interior engine}:        \> \emph{special}\\
    \refKey{/tcb/segmentation engine}:    \> \emph{special}\\
    \refKey{/tcb/title engine}:           \> |pathmiddle|
  \end{tabbing}
\end{tcolorbox}
\end{docSkin}


\begin{dispExample}
\skinExampleSet{skin=bicolormiddle,
  colbacklower=LimeGreen!75!LightGreen,
}
\end{dispExample}


\clearpage
\begin{docSkin}{bicolorlast}
这是 \refSkin{bicolor} 的一个变体，用作 \refSkin{bicolor} 断开序列的\emph{末}段。
不过，该皮肤也可以独立使用。
\begin{tcolorbox}[skintable=bicolorlast]
  \begin{tabbing}
    \refKey{/tcb/interior titled engine}: \=\kill
    \refKey{/tcb/frame engine}:           \> |pathlast|\\
    \refKey{/tcb/interior titled engine}: \> \emph{special}\\ 
    \refKey{/tcb/interior engine}:        \> \emph{special}\\
    \refKey{/tcb/segmentation engine}:    \> \emph{special}\\
    \refKey{/tcb/title engine}:           \> |pathlast|
  \end{tabbing}
\end{tcolorbox}
\end{docSkin}


\begin{dispExample}
\skinExampleSet{skin=bicolorlast,
  colbacklower=LimeGreen!75!LightGreen,
}
\end{dispExample}


\clearpage

\begin{docSkin}[doc new=2021-05-21]{bicolor jigsaw}
  这是 \refSkin{bicolor} 皮肤的拼图变体。区别在于其边框引擎，
  见 \zvref{subsec:skinjigsaw}。
\begin{tcolorbox}[skintable=bicolor jigsaw]
  \begin{tabbing}
    \refKey{/tcb/interior titled engine}: \=\kill
    \refKey{/tcb/frame engine}:           \> |pathjigsaw|\\
    \refKey{/tcb/interior titled engine}: \> \emph{special}\\ 
    \refKey{/tcb/interior engine}:        \> \emph{special}\\
    \refKey{/tcb/segmentation engine}:    \> \emph{special}\\
    \refKey{/tcb/title engine}:           \> |path|
  \end{tabbing}
\end{tcolorbox}
\end{docSkin}

\begin{docTcbKey}{bicolor jigsaw}{}{样式，无取值}
  这是设置 |skin=enhanced jigsaw| 的简写。
\end{docTcbKey}


\begin{dispExample*}{segmentation style={preaction={fill=white},pattern=checkerboard,pattern color=gray!40}}
\skinExampleSet{bicolor jigsaw,
  colbacklower=LimeGreen!75!LightGreen,
  opacityframe=0.5,opacityback=0.5,opacitybacktitle=0.5,
  opacitybacklower=0.5,
}
\end{dispExample*}


\clearpage


\begin{docSkin}[doc new=2021-05-21]{bicolorfirst jigsaw}
  这是 \refSkin{bicolorfirst} 皮肤的拼图变体。区别在于其边框引擎，
  见 \zvref{subsec:skinjigsaw}。
\begin{tcolorbox}[skintable=bicolorfirst jigsaw]
  \begin{tabbing}
    \refKey{/tcb/interior titled engine}: \=\kill
    \refKey{/tcb/frame engine}:           \> |pathfirstjigsaw|\\
    \refKey{/tcb/interior titled engine}: \> \emph{special}\\ 
    \refKey{/tcb/interior engine}:        \> \emph{special}\\
    \refKey{/tcb/segmentation engine}:    \> \emph{special}\\
    \refKey{/tcb/title engine}:           \> |pathfirst|
  \end{tabbing}
\end{tcolorbox}
\end{docSkin}

\begin{dispExample*}{segmentation style={preaction={fill=white},pattern=checkerboard,pattern color=gray!40}}
\skinExampleSet{skin=bicolorfirst jigsaw,
  colbacklower=LimeGreen!75!LightGreen,
  opacityframe=0.5,opacityback=0.5,opacitybacktitle=0.5,
  opacitybacklower=0.5,
}
\end{dispExample*}



\clearpage

\begin{docSkin}[doc new=2021-05-21]{bicolormiddle jigsaw}
  这是 \refSkin{bicolormiddle} 皮肤的拼图变体。区别在于其边框引擎，
  见 \zvref{subsec:skinjigsaw}。
\begin{tcolorbox}[skintable=bicolormiddle jigsaw]
  \begin{tabbing}
    \refKey{/tcb/interior titled engine}: \=\kill
    \refKey{/tcb/frame engine}:           \> |pathmiddlejigsaw|\\
    \refKey{/tcb/interior titled engine}: \> \emph{special}\\ 
    \refKey{/tcb/interior engine}:        \> \emph{special}\\
    \refKey{/tcb/segmentation engine}:    \> \emph{special}\\
    \refKey{/tcb/title engine}:           \> |pathmiddle|
  \end{tabbing}
\end{tcolorbox}
\end{docSkin}


\begin{dispExample*}{segmentation style={preaction={fill=white},pattern=checkerboard,pattern color=gray!40}}
\skinExampleSet{skin=bicolormiddle jigsaw,
  colbacklower=LimeGreen!75!LightGreen,
  opacityframe=0.5,opacityback=0.5,opacitybacktitle=0.5,
  opacitybacklower=0.5,
}
\end{dispExample*}


\clearpage
\begin{docSkin}[doc new=2021-05-21]{bicolorlast jigsaw}
  这是 \refSkin{bicolorlast} 皮肤的拼图变体。区别在于其边框引擎，
  见 \zvref{subsec:skinjigsaw}。
\begin{tcolorbox}[skintable=bicolorlast jigsaw]
  \begin{tabbing}
    \refKey{/tcb/interior titled engine}: \=\kill
    \refKey{/tcb/frame engine}:           \> |pathlastjigsaw|\\
    \refKey{/tcb/interior titled engine}: \> \emph{special}\\ 
    \refKey{/tcb/interior engine}:        \> \emph{special}\\
    \refKey{/tcb/segmentation engine}:    \> \emph{special}\\
    \refKey{/tcb/title engine}:           \> |pathlast|
  \end{tabbing}
\end{tcolorbox}
\end{docSkin}


\begin{dispExample*}{segmentation style={preaction={fill=white},pattern=checkerboard,pattern color=gray!40}}
\skinExampleSet{skin=bicolorlast jigsaw,
  colbacklower=LimeGreen!75!LightGreen,
  opacityframe=0.5,opacityback=0.5,opacitybacktitle=0.5,
  opacitybacklower=0.5,
}
\end{dispExample*}



\clearpage
\subsection{皮肤家族 \enquote{tile}}
\begin{docSkin}[doc new=2016-02-25]{tile}
  该皮肤是 \refSkin{bicolor} 皮肤的一个变体。特别地，文本框的可选下部
  由 \refKey{/tcb/colbacklower} 着色。与 \refSkin{bicolor} 的主要区别在于
  \refSkin{tile} 没有边框。
\begin{tcolorbox}[skintable=tile]
  \begin{tabbing}
    \refKey{/tcb/interior titled engine}: \=\kill
    \refKey{/tcb/frame engine}:           \> |empty|\\
    \refKey{/tcb/interior titled engine}: \> \emph{special}\\ 
    \refKey{/tcb/interior engine}:        \> \emph{special}\\
    \refKey{/tcb/segmentation engine}:    \> \emph{special}\\
    \refKey{/tcb/title engine}:           \> |path|
  \end{tabbing}
\end{tcolorbox}
\end{docSkin}

\begin{docTcbKey}[][doc new=2016-02-25]{tile}{}{样式，初始\\
  |skin=tile,boxrule=0pt,sharp corners,title filled,fonttitle=\textbackslash bfseries|
}
  本键应用 |skin=tile|，此外还会改变几何尺寸和一些样式选项。
\end{docTcbKey}


\begin{dispExample}
\skinExampleSet{tile,
  colbacklower=LimeGreen!75!LightGreen,
}
\end{dispExample}


\clearpage
\begin{docSkin}[doc new=2016-02-25]{tilefirst}
这是 \refSkin{tile} 的一个变体，用作 \refSkin{tile} 断开序列的\emph{首}段。
不过，该皮肤也可以独立使用。
\begin{tcolorbox}[skintable=tilefirst]
  \begin{tabbing}
    \refKey{/tcb/interior titled engine}: \=\kill
    \refKey{/tcb/frame engine}:           \> |empty|\\
    \refKey{/tcb/interior titled engine}: \> \emph{special}\\ 
    \refKey{/tcb/interior engine}:        \> \emph{special}\\
    \refKey{/tcb/segmentation engine}:    \> \emph{special}\\
    \refKey{/tcb/title engine}:           \> |pathfirst|
  \end{tabbing}
\end{tcolorbox}
\end{docSkin}

\begin{dispExample}
\skinExampleSet{skin=tilefirst,
  colbacklower=LimeGreen!75!LightGreen,
  boxrule=0pt,
}
\end{dispExample}


\clearpage
\begin{docSkin}[doc new=2016-02-25]{tilemiddle}
这是 \refSkin{tile} 的一个变体，用作 \refSkin{tile} 断开序列的\emph{中间}段。
不过，该皮肤也可以独立使用。
\begin{tcolorbox}[skintable=tilemiddle]
  \begin{tabbing}
    \refKey{/tcb/interior titled engine}: \=\kill
    \refKey{/tcb/frame engine}:           \> |empty|\\
    \refKey{/tcb/interior titled engine}: \> \emph{special}\\ 
    \refKey{/tcb/interior engine}:        \> \emph{special}\\
    \refKey{/tcb/segmentation engine}:    \> \emph{special}\\
    \refKey{/tcb/title engine}:           \> |pathmiddle|
  \end{tabbing}
\end{tcolorbox}
\end{docSkin}


\begin{dispExample}
\skinExampleSet{skin=tilemiddle,
  colbacklower=LimeGreen!75!LightGreen,
  boxrule=0pt,
}
\end{dispExample}


\clearpage
\begin{docSkin}[doc new=2016-02-25]{tilelast}
这是 \refSkin{tile} 的一个变体，用作 \refSkin{tile} 断开序列的\emph{末}段。
不过，该皮肤也可以独立使用。
\begin{tcolorbox}[skintable=tilelast]
  \begin{tabbing}
    \refKey{/tcb/interior titled engine}: \=\kill
    \refKey{/tcb/frame engine}:           \> |empty|\\
    \refKey{/tcb/interior titled engine}: \> \emph{special}\\ 
    \refKey{/tcb/interior engine}:        \> \emph{special}\\
    \refKey{/tcb/segmentation engine}:    \> \emph{special}\\
    \refKey{/tcb/title engine}:           \> |pathlast|
  \end{tabbing}
\end{tcolorbox}
\end{docSkin}


\begin{dispExample}
\skinExampleSet{skin=tilelast,
  colbacklower=LimeGreen!75!LightGreen,
  boxrule=0pt,
}
\end{dispExample}



\clearpage
\subsection{皮肤家族 \enquote{beamer}}

\begin{docSkin}{beamer}
  该皮肤模仿 \refPkg{beamer} 文档类中常见的文本框，因此称为 \enquote{beamer}。
  它使用核心宏包的常规颜色，但略微加深了一点。\par
  设计该皮肤的初衷是在 \refPkg{beamer} 文档类之外使用具有\enquote{beamer风格}的文本框；
  不过，你也可能想在 |beamer| 幻灯片中用 |tcolorbox| 作为块。为此可参见
  \refPkg{beamertheme-tcolorbox} 宏包，它提供了一个 |beamer| 内部主题，
  用 |tcolorbox| 重现标准的 |beamer| 块。
\begin{tcolorbox}[skintable=beamer]
  \begin{tabbing}
    \refKey{/tcb/interior titled engine}: \=\kill
    \refKey{/tcb/frame engine}:           \> |path|\\
    \refKey{/tcb/interior titled engine}: \> \emph{special}\\ 
    \refKey{/tcb/interior engine}:        \> \emph{special}\\
    \refKey{/tcb/segmentation engine}:    \> \emph{special}\\
    \refKey{/tcb/title engine}:           \> |path|
  \end{tabbing}
\end{tcolorbox}
\end{docSkin}



\begin{docTcbKey}{beamer}{}{样式，初始\\
  |skin=beamer,boxrule=0mm,titlerule=1mm,toptitle=0.5mm,|\\
  |arc=2mm,fonttitle=\textbackslash bfseries,drop fuzzy shadow|
}
  本键应用 |skin=beamer|，此外还会改变几何尺寸和一些样式选项。
\end{docTcbKey}



\begin{dispExample}
\skinExampleSet{beamer,title filled=false}
\end{dispExample}



\begin{dispExample}
\begin{tcolorbox}[beamer,colback=Salmon!50!white,colframe=FireBrick!75!black,
  adjusted title=A colored box with the \enquote{beamer} skin]
This box looks like a box provided by the \texttt{beamer} class.
\end{tcolorbox}
\end{dispExample}


\begin{dispExample}
\begin{tcolorbox}[beamer,colframe=blue,colback=black,
  watermark graphics=lichtspiel.jpg,
  coltext=white,watermark opacity=0.75,watermark stretch=1.0,
  title=Beamer Box with background picture]
\lipsum[1]
\end{tcolorbox}
\end{dispExample}


\begin{dispExample}
\newtcolorbox{myblock}[2][]{%
  beamer,breakable,colback=LightBlue,colframe=DarkBlue,#1,title={#2}}%

\begin{myblock}{Beamerish \texttt{block}: \texttt{myblock}}
\lipsum[1]
\end{myblock}
\end{dispExample}


\clearpage
\begin{docSkin}{beamerfirst}
这是 \refSkin{beamer} 的一个变体，用作 \refSkin{beamer} 断开序列的\emph{首}段。
不过，该皮肤也可以独立使用。
\begin{tcolorbox}[skintable=beamerfirst]
  \begin{tabbing}
    \refKey{/tcb/interior titled engine}: \=\kill
    \refKey{/tcb/frame engine}:           \> |pathfirst|\\
    \refKey{/tcb/interior titled engine}: \> \emph{special}\\ 
    \refKey{/tcb/interior engine}:        \> \emph{special}\\
    \refKey{/tcb/segmentation engine}:    \> \emph{special}\\
    \refKey{/tcb/title engine}:           \> |pathfirst|
  \end{tabbing}
\end{tcolorbox}
\end{docSkin}


\begin{dispExample}
\skinExampleSet{beamer,title filled=false,skin=beamerfirst}
\end{dispExample}


\medskip

\begin{docSkin}{beamermiddle}
这是 \refSkin{beamer} 的一个变体，用作 \refSkin{beamer} 断开序列的\emph{中间}段。
不过，该皮肤也可以独立使用。
\begin{tcolorbox}[skintable=beamermiddle]
  \begin{tabbing}
    \refKey{/tcb/interior titled engine}: \=\kill
    \refKey{/tcb/frame engine}:           \> |pathmiddle|\\
    \refKey{/tcb/interior titled engine}: \> \emph{special}\\ 
    \refKey{/tcb/interior engine}:        \> \emph{special}\\
    \refKey{/tcb/segmentation engine}:    \> \emph{special}\\
    \refKey{/tcb/title engine}:           \> |pathmiddle|
  \end{tabbing}
\end{tcolorbox}
\end{docSkin}


\begin{dispExample}
\skinExampleSet{beamer,title filled=false,skin=beamermiddle}
\end{dispExample}


\clearpage
\begin{docSkin}{beamerlast}
这是 \refSkin{beamer} 的一个变体，用作 \refSkin{beamer} 断开序列的\emph{末}段。
不过，该皮肤也可以独立使用。
\begin{tcolorbox}[skintable=beamerlast]
  \begin{tabbing}
    \refKey{/tcb/interior titled engine}: \=\kill
    \refKey{/tcb/frame engine}:           \> |pathlast|\\
    \refKey{/tcb/interior titled engine}: \> \emph{special}\\ 
    \refKey{/tcb/interior engine}:        \> \emph{special}\\
    \refKey{/tcb/segmentation engine}:    \> \emph{special}\\
    \refKey{/tcb/title engine}:           \> |pathlast|
  \end{tabbing}
\end{tcolorbox}
\end{docSkin}

\begin{dispExample}
\skinExampleSet{beamer,title filled=false,skin=beamerlast}
\end{dispExample}



\clearpage
\subsection{皮肤家族 \enquote{widget}}
\begin{docSkin}{widget}
  该皮肤使用核心宏包的常规颜色，但略微加深了一点。
  如有需要，可通过 \refKey{/tcb/frame style}、
  \refKey{/tcb/interior style} 和 \refKey{/tcb/segmentation style}
  控制该皮肤的外观。
\begin{tcolorbox}[skintable=widget]
  \begin{tabbing}
    \refKey{/tcb/interior titled engine}: \=\kill
    \refKey{/tcb/frame engine}:           \> |path|\\
    \refKey{/tcb/interior titled engine}: \> |path|\\ 
    \refKey{/tcb/interior engine}:        \> |path|\\
    \refKey{/tcb/segmentation engine}:    \> \emph{special}\\
    \refKey{/tcb/title engine}:           \> \emph{special}
  \end{tabbing}
\end{tcolorbox}
\end{docSkin}


\begin{docTcbKey}[][doc updated={2020-09-23}]{widget}{}{样式，初始\\
  |skin=widget,arc=0.5mm,fonttitle=\bfseries,titlerule=0mm|
}
  本键应用 |skin=widget|，此外还会改变几何尺寸和一些样式选项。
\end{docTcbKey}


\begin{dispExample}
\skinExampleSet{widget}
\end{dispExample}


\begin{dispExample}
\begin{tcolorbox}[widget,colback=Salmon!50!white,colframe=FireBrick!75!black,
  adjusted title=A colored box with the \enquote{widget} skin]
This is my content.
\end{tcolorbox}
\end{dispExample}

\clearpage

\begin{docSkin}{widgetfirst}
这是 \refSkin{widget} 的一个变体，用作 \refSkin{widget} 断开序列的\emph{首}段。
不过，该皮肤也可以独立使用。
\begin{tcolorbox}[skintable=widgetfirst]
  \begin{tabbing}
    \refKey{/tcb/interior titled engine}: \=\kill
    \refKey{/tcb/frame engine}:           \> |pathfirst|\\
    \refKey{/tcb/interior titled engine}: \> |pathfirst|\\ 
    \refKey{/tcb/interior engine}:        \> |pathfirst|\\
    \refKey{/tcb/segmentation engine}:    \> \emph{special}\\
    \refKey{/tcb/title engine}:           \> \emph{special}
  \end{tabbing}
\end{tcolorbox}
\end{docSkin}


\begin{dispExample}
\skinExampleSet{widget,skin=widgetfirst}
\end{dispExample}

\medskip

\begin{docSkin}{widgetmiddle}
这是 \refSkin{widget} 的一个变体，用作 \refSkin{widget} 断开序列的\emph{中间}段。
不过，该皮肤也可以独立使用。
\begin{tcolorbox}[skintable=widgetmiddle]
  \begin{tabbing}
    \refKey{/tcb/interior titled engine}: \=\kill
    \refKey{/tcb/frame engine}:           \> |pathmiddle|\\
    \refKey{/tcb/interior titled engine}: \> |pathmiddle|\\ 
    \refKey{/tcb/interior engine}:        \> |pathmiddle|\\
    \refKey{/tcb/segmentation engine}:    \> \emph{special}\\
    \refKey{/tcb/title engine}:           \> \emph{special}
  \end{tabbing}
\end{tcolorbox}
\end{docSkin}

\begin{dispExample}
\skinExampleSet{widget,skin=widgetmiddle}
\end{dispExample}


\clearpage
\begin{docSkin}{widgetlast}
这是 \refSkin{widget} 的一个变体，用作 \refSkin{widget} 断开序列的\emph{末}段。
不过，该皮肤也可以独立使用。
\begin{tcolorbox}[skintable=widgetlast]
  \begin{tabbing}
    \refKey{/tcb/interior titled engine}: \=\kill
    \refKey{/tcb/frame engine}:           \> |pathlast|\\
    \refKey{/tcb/interior titled engine}: \> |pathlast|\\ 
    \refKey{/tcb/interior engine}:        \> |pathlast|\\
    \refKey{/tcb/segmentation engine}:    \> \emph{special}\\
    \refKey{/tcb/title engine}:           \> \emph{special}
  \end{tabbing}
\end{tcolorbox}
\end{docSkin}


\begin{dispExample}
\skinExampleSet{widget,skin=widgetlast}
\end{dispExample}


\clearpage
\subsection{皮肤家族 \enquote{empty}}

\begin{docSkin}{empty}
  该皮肤把所有引擎都设为 |empty|，即完全不绘制任何东西。
  因此，该皮肤是你自己从头创建全新样式的良好起点。
\begin{tcolorbox}[skintable=empty]
  \begin{tabbing}
    \refKey{/tcb/interior titled engine}: \=\kill
    \refKey{/tcb/frame engine}:           \> |empty|\\
    \refKey{/tcb/interior titled engine}: \> |empty|\\ 
    \refKey{/tcb/interior engine}:        \> |empty|\\
    \refKey{/tcb/segmentation engine}:    \> |empty|\\
    \refKey{/tcb/title engine}:           \> |empty|
  \end{tabbing}
\end{tcolorbox}
\end{docSkin}


\begin{marker}
  请注意，应用皮肤时文字颜色保持不变。
  由于默认标题颜色是白色，若不用 \refKey{/tcb/coltitle} 设为其他颜色，
  采用 \docValue*{empty} 皮肤的文本框标题会变得不可见。
\end{marker}


\begin{docTcbKey}{empty}{}{样式，无取值}
  这是设置 |skin=empty| 的简写。
\end{docTcbKey}


\begin{dispExample}
\skinExampleSet{empty,
  coltitle=Navy,borderline={2pt}{0pt}{black!10!white},
}
\end{dispExample}


\clearpage
\begin{docTcbKey}{blanker}{}{样式，初始未设置}
  本样式基于 \refSkin{empty} 皮肤。所有引擎均设为 empty，
  所有边距均设为 |0pt|。与 \refKey{/tcb/blank} 不同，除几何节点外，
  图形路径不会被构建。
\begin{dispExample}
\begin{tcolorbox}[blanker,watermark text=A blank box]
\lipsum[1]
\end{tcolorbox}
\end{dispExample}

\begin{dispExample}
% \tcbuselibrary{fitting}
\newtcboxfit{\mybox}[1]{blanker,width=4cm,height=7cm,top=4pt,
  watermark text=#1}

\begin{tabular}{|c|c|c|}\hline
A & B & C\\\hline
\mybox{A}{\lipsum[1]} & \mybox{B}{\lipsum[2]} & \mybox{C}{\lipsum[3]}\\\hline
\end{tabular}
\end{dispExample}
\end{docTcbKey}



\clearpage
\begin{docTcbKey}{blankest}{}{样式，初始未设置}
  本样式在 \refKey{/tcb/blanker} 的基础上扩展。
  所有引擎均设为 empty，所有边距均设为 |0pt|。
  与 \refKey{/tcb/blanker} 不同，标题、阴影、底层装饰、
  覆盖层、收尾装饰和边线也都被移除。

\begin{dispExample}
% \tcbuselibrary{raster}
\begin{tcbraster}[raster columns=3,raster equal height,
    title=Box \thetcbrasternum,
    enhanced,size=small,colframe=red!50!black,colback=red!10!white,
    coltitle=yellow!85!black,
    drop fuzzy shadow,watermark text={Box \thetcbrasternum},
    borderline={.25mm}{-0.5mm}{green!40!black},
    finish={\begin{tcbclipframe}\draw[blue,opacity=0.1,line width=1cm]
      (frame.south west) -- (frame.north east);\end{tcbclipframe}},
    ]
  \begin{tcolorbox}\lipsum[4]\end{tcolorbox}
  \begin{tcolorbox}[blanker]\lipsum[4]\end{tcolorbox}
  \begin{tcolorbox}[blankest]\lipsum[4]\end{tcolorbox}
\end{tcbraster}
\end{dispExample}
\end{docTcbKey}


\clearpage
\begin{docSkin}{emptyfirst}
这是 \refSkin{empty} 的一个变体，用作 \refSkin{empty} 断开序列的\emph{首}段。
不过，该皮肤也可以独立使用。
\begin{tcolorbox}[skintable=emptyfirst]
  \begin{tabbing}
    \refKey{/tcb/interior titled engine}: \=\kill
    \refKey{/tcb/frame engine}:           \> |empty|\\
    \refKey{/tcb/interior titled engine}: \> |empty|\\ 
    \refKey{/tcb/interior engine}:        \> |empty|\\
    \refKey{/tcb/segmentation engine}:    \> |empty|\\
    \refKey{/tcb/title engine}:           \> |empty|
  \end{tabbing}
\end{tcolorbox}
\end{docSkin}


\begin{dispExample}
\skinExampleSet{skin=emptyfirst,
  coltitle=Navy,borderline={2pt}{0pt}{black!10!white},
}
\end{dispExample}


\clearpage

\begin{docSkin}{emptymiddle}
这是 \refSkin{empty} 的一个变体，用作 \refSkin{empty} 断开序列的\emph{中间}段。
不过，该皮肤也可以独立使用。
\begin{tcolorbox}[skintable=emptymiddle]
  \begin{tabbing}
    \refKey{/tcb/interior titled engine}: \=\kill
    \refKey{/tcb/frame engine}:           \> |empty|\\
    \refKey{/tcb/interior titled engine}: \> |empty|\\ 
    \refKey{/tcb/interior engine}:        \> |empty|\\
    \refKey{/tcb/segmentation engine}:    \> |empty|\\
    \refKey{/tcb/title engine}:           \> |empty|
  \end{tabbing}
\end{tcolorbox}
\end{docSkin}


\begin{dispExample}
\skinExampleSet{skin=emptymiddle,
  coltitle=Navy,borderline={2pt}{0pt}{black!10!white},
}
\end{dispExample}


\clearpage
\begin{docSkin}{emptylast}
这是 \refSkin{empty} 的一个变体，用作 \refSkin{empty} 断开序列的\emph{末}段。
不过，该皮肤也可以独立使用。
\begin{tcolorbox}[skintable=emptylast]
  \begin{tabbing}
    \refKey{/tcb/interior titled engine}: \=\kill
    \refKey{/tcb/frame engine}:           \> |empty|\\
    \refKey{/tcb/interior titled engine}: \> |empty|\\ 
    \refKey{/tcb/interior engine}:        \> |empty|\\
    \refKey{/tcb/segmentation engine}:    \> |empty|\\
    \refKey{/tcb/title engine}:           \> |empty|
  \end{tabbing}
\end{tcolorbox}
\end{docSkin}

\begin{dispExample}
\skinExampleSet{skin=emptylast,
  coltitle=Navy,borderline={2pt}{0pt}{black!10!white},
}
\end{dispExample}

\clearpage
\begin{dispListing*}{breakable,phantomlabel=freeboxexample,before upper={本示例展示
一个可断开的自定义文本框。这里我们定义一个环境 |freebox|。
第一次使用 |freebox| 产生一个未断开的 |tcolorbox|。
该文本框由 \refKey{/tcb/frame code}
和 \refKey{/tcb/interior code} 给出的代码绘制。\par
第二次使用 |freebox| 时，它被断成若干段，
各段由 \refKey{/tcb/skin first is subskin of}、
\refKey{/tcb/skin middle is subskin of} 和
\refKey{/tcb/skin last is subskin of} 给出的代码绘制。
\par\bigskip
}}
% Preamble:
%\usepackage{lipsum}
%\tcbuselibrary{skins,breakable}
\tikzset{coltria/.style={fill=red!15!white}}

\newtcolorbox{freebox}[1][]{empty,
  breakable,height fixed for=first and middle,
  leftrule=5mm,left=2mm,
  frame style={fill,top color=red!75!black,bottom color=red!75!black,middle color=red},
  colback=yellow!50!white,
  watermark color=red!50!yellow!75!white,
  watermark text on=unbroken is unbroken box,
  watermark text on=first is first part,
  watermark text on=middle is middle part,
  watermark text on=last is last part,
  % code for unbroken boxes:
  frame code={\path[tcb fill frame] (frame.south west)--(frame.north west)
    --([xshift=-5mm]frame.north east)--([yshift=-5mm]frame.north east)
    --([yshift=5mm]frame.south east)--([xshift=-5mm]frame.south east)--cycle; },
  interior code={\path[tcb fill interior] (interior.south west)--(interior.north west)
    --([xshift=-4.8mm]interior.north east)--([yshift=-4.8mm]interior.north east)
    --([yshift=4.8mm]interior.south east)--([xshift=-4.8mm]interior.south east)
    --cycle; },
  % code for the first part of a break sequence:
  skin first is subskin of={emptyfirst}{%
    frame code={\path[tcb fill frame] (frame.south west)--(frame.north west)
      --([xshift=-5mm]frame.north east)--([yshift=-5mm]frame.north east)
      --(frame.south east)--cycle;
      \path[coltria] ([xshift=2.5mm,yshift=1mm]frame.south west) -- +(120:2mm)
      -- +(60:2mm)-- cycle; },
    interior code={\path[tcb fill interior] (interior.south west|-frame.south)
      --(interior.north west)--([xshift=-4.8mm]interior.north east)
      --([yshift=-4.8mm]interior.north east)--(interior.south east|-frame.south)
      --cycle; },
  },%
  % code for the middle part of a break sequence:
  skin middle is subskin of={emptymiddle}{%
    frame code={\path[tcb fill frame] (frame.south west)--(frame.north west)
      --(frame.north east)--(frame.south east)--cycle;
      \path[coltria] ([xshift=2.5mm,yshift=-1mm]frame.north west) -- +(240:2mm)
        -- +(300:2mm) -- cycle;
      \path[coltria] ([xshift=2.5mm,yshift=1mm]frame.south west) -- +(120:2mm)
        -- +(60:2mm) -- cycle;
      },
    interior code={\path[tcb fill interior] (interior.south west|-frame.south)
      --(interior.north west|-frame.north)--(interior.north east|-frame.north)
      --(interior.south east|-frame.south)--cycle; },
    },
  % code for the last part of a break sequence:
  skin last is subskin of={emptylast}{%
    frame code={\path[tcb fill frame] (frame.south west)--(frame.north west)
      --(frame.north east)--([yshift=5mm]frame.south east)
      --([xshift=-5mm]frame.south east)--cycle;
      \path[coltria] ([xshift=2.5mm,yshift=-1mm]frame.north west) -- +(240:2mm)
      -- +(300:2mm) -- cycle;
      },
    interior code={\path[tcb fill interior] (interior.south west)
      --(interior.north west|-frame.north)--(interior.north east|-frame.north)
      --([yshift=4.8mm]interior.south east)--([xshift=-4.8mm]interior.south east)
      --cycle; },
    },
  #1}

\begin{freebox}
\lipsum[1]
\end{freebox}

\begin{freebox}
\lipsum[1-12]
\end{freebox}
\end{dispListing*}
{\tcbusetemp}


\clearpage

\subsection{皮肤 \enquote{spartan}}\label{subsec:spartan}

\begin{docSkin}{spartan}
  该皮肤相当\ldots\,简朴。它不支持圆角、覆盖层、阴影、边线和收尾装饰。
  唯一的例外是底层装饰。用这种皮肤做不出很花哨的效果，但编译非常快。
  因此，|spartan| 皮肤被用于草稿模式，见 \zvref{subsec:draftmode}。
  不过，它也可以当作普通皮肤使用。

\begin{tcolorbox}[skintable=spartan]
  \begin{tabbing}
    \refKey{/tcb/interior titled engine}: \=\kill
    \refKey{/tcb/frame engine}:           \> |spartan|\\
    \refKey{/tcb/interior titled engine}: \> |spartan|\\ 
    \refKey{/tcb/interior engine}:        \> |spartan|\\
    \refKey{/tcb/segmentation engine}:    \> |spartan|\\
    \refKey{/tcb/title engine}:           \> |spartan|
  \end{tabbing}
\end{tcolorbox}
\end{docSkin}


\begin{docTcbKey}{spartan}{}{样式，无取值}
  这是设置 |skin=spartan| 的简写。
\end{docTcbKey}


\begin{dispExample}
\skinExampleSet{spartan}
\end{dispExample}


\clearpage

\subsection{皮肤 \enquote{draft}}\label{subsec:draft}

\begin{docSkin}{draft}
  该皮肤用于在为 |tcolorbox| 起草新的几何设置时使用。
\begin{tcolorbox}[skintable=draft]
  \begin{tabbing}
    \refKey{/tcb/interior titled engine}: \=\kill
    \refKey{/tcb/frame engine}:           \> \emph{special}\\
    \refKey{/tcb/interior titled engine}: \> \emph{special}\\ 
    \refKey{/tcb/interior engine}:        \> \emph{special}\\
    \refKey{/tcb/segmentation engine}:    \> |path|\\
    \refKey{/tcb/title engine}:           \> |path|
  \end{tabbing}
\end{tcolorbox}
\end{docSkin}

\begin{docTcbKey}{draft}{}{样式，无取值}
  这是设置 |skin=draft| 的简写。
\end{docTcbKey}


\begin{dispExample}
\skinExampleSet{draft}
\end{dispExample}



\begin{dispExample}
\vspace*{3mm}
\begin{tcolorbox}[draft,title=A colored box with the \enquote{draft} skin]
\lipsum[1-3]
\tcblower
\lipsum[4-6]
\end{tcolorbox}
\end{dispExample}



\clearpage
\subsection{皮肤家族 \enquote{freelance}}
\begin{marker}
这个 \enquote{freelance} 皮肤家族自 |tcolorbox| 3.00 起已弃用。
它不再是必需的，因为 \refKey{/tcb/frame code}、
\refKey{/tcb/interior code}、
\refKey{/tcb/interior titled code} 和
\refKey{/tcb/title code} 现在可应用于任何皮肤。从这个意义上说，
现在一切都已变成\emph{自由}的了。\par
对于 \refKey{/tcb/freelance} 的用户：旧代码应仍能继续工作。
在某些情况下，可断开的 freelance 文本框可能会有例外。
对于新代码，在你过去会使用 \refKey{/tcb/freelance} 的地方，
请改用 \refKey{/tcb/empty} 或 \refKey{/tcb/enhanced}。
\end{marker}

\begin{docSkin}{freelance}
  该皮肤让 |tcolorbox| 的外观完全自由。
  所有绘制引擎均设为 |freelance| 类型；它们使用 \refPkg{tikz} 宏包，
  并计算 \refKey{/tcb/geometry nodes}。
  %This skin is useful for boxes which should differ much from the normal
  %appearance. Note that this difference has to be programmed by the user.
  %The drawing code can be given
  %with the following option keys. As default value, the code from the |standard|
  %skin is set.
\begin{tcolorbox}[skintable=freelance]
  \begin{tabbing}
    \refKey{/tcb/interior titled engine}: \=\kill
    \refKey{/tcb/frame engine}:           \> |freelance|\\
    \refKey{/tcb/interior titled engine}: \> |freelance|\\ 
    \refKey{/tcb/interior engine}:        \> |freelance|\\
    \refKey{/tcb/segmentation engine}:    \> |freelance|\\
    \refKey{/tcb/title engine}:           \> |freelance|
  \end{tabbing}
\end{tcolorbox}
\end{docSkin}

\begin{docTcbKey}{freelance}{}{样式，无取值}
  这是设置 |skin=freelance| 的简写。
\end{docTcbKey}

\begin{docSkin}{freelancefirst}
  该皮肤与 \refSkin{freelance} 相同，只是断开序列不同，
  见 \zvref{subsec:breaksequence}。
  %It is used as first part of the
  %break sequence of \refSkin{freelance}. \refKey{/tcb/extend freelancefirst}
  %can be used to customize this part.
\end{docSkin}

\begin{docSkin}{freelancemiddle}
  该皮肤与 \refSkin{freelance} 相同，只是断开序列不同，
  见 \zvref{subsec:breaksequence}。
  %It is used as middle part of the
  %break sequence of \refSkin{freelance}. \refKey{/tcb/extend freelancemiddle}
  %can be used to customize this part.
\end{docSkin}

\begin{docSkin}{freelancelast}
  该皮肤与 \refSkin{freelance} 相同，只是断开序列不同，
  见 \zvref{subsec:breaksequence}。
  %It is used as last part of the
  %break sequence of \refSkin{freelance}. \refKey{/tcb/extend freelancelast}
  %can be used to customize this part.
\end{docSkin}


\begin{docTcbKey}{extend freelance}{=\meta{options}}{无默认值，初始为空}
把 \meta{选项} 添加到 \refSkin{freelance} 的皮肤定义中。
\end{docTcbKey}

\begin{docTcbKey}{extend freelancefirst}{=\meta{options}}{无默认值，初始为空}
把 \meta{选项} 添加到 \refSkin{freelancefirst} 的皮肤定义中，
后者用作 \refSkin{freelance} 断开序列的首段。本键的替代选项见
\refKey{/tcb/skin first is subskin of}。
\end{docTcbKey}

\begin{docTcbKey}{extend freelancemiddle}{=\meta{options}}{无默认值，初始为空}
把 \meta{选项} 添加到 \refSkin{freelancemiddle} 的皮肤定义中，
后者用作 \refSkin{freelance} 断开序列的中间段。本键的替代选项见
\refKey{/tcb/skin middle is subskin of}。
\end{docTcbKey}

\enlargethispage*{1cm}

\begin{docTcbKey}{extend freelancelast}{=\meta{options}}{无默认值，初始为空}
把 \meta{选项} 添加到 \refSkin{freelancelast} 的皮肤定义中，
后者用作 \refSkin{freelance} 断开序列的末段。本键的替代选项见
\refKey{/tcb/skin last is subskin of}。
\end{docTcbKey}



